\subsection{Tối ưu hiệu năng hệ thống}

Một số kỹ thuật tối ưu hiệu năng được áp dụng nhằm giữ độ trễ ở mức chấp nhận được dưới tải bầu cử đỉnh -- khi hàng nghìn cử tri có thể đồng thời bỏ phiếu trong cùng một khoảng thời gian ngắn.

\textbf{MongoDB Replica Set và Prisma transaction:} Cụm MongoDB được khởi tạo với \pkg{--replSet rs0}, là điều kiện bắt buộc để Prisma sử dụng \pkg{\$transaction} với nhiều thao tác ACID. Hệ thống tuân thủ \textit{kiểm tra hợp lệ ngoài giao dịch, chỉ ghi (commit) trong giao dịch}: các thao tác đọc và side-effect nặng (gọi Fabric mất 2--3 giây) được thực hiện \textbf{ngoài} giao dịch; giao dịch chỉ chứa hai bước \textit{re-validate} (chống xung đột đồng thời) và \textit{ghi}. Nhờ đó thời gian giữ khóa của giao dịch ngắn, thông lượng tăng đáng kể so với cách bao bọc toàn bộ luồng nghiệp vụ.

\textbf{Cache phản hồi HTTP:} \pkg{HttpCacheInterceptor} đệm phản hồi của các yêu cầu \pkg{GET} vào Redis. Các endpoint danh sách bầu cử công khai và xem chi tiết ứng viên hưởng lợi nhiều nhất từ cơ chế này vì chúng có tải đọc cao và dữ liệu thay đổi ít.

\textbf{Chiến lược chỉ mục MongoDB:} Các bảng chính được phủ chỉ mục để mọi thao tác đọc theo dõi được $O(\log N)$ hoặc $O(1)$ thay vì quét tuần tự. Bảng~\ref{tab:mongo-index} liệt kê các chỉ mục then chốt.

\begin{table}[H]
\centering
\caption{Chiến lược chỉ mục các bảng chính trong MongoDB}
\label{tab:mongo-index}
\begin{tabularx}{\textwidth}{|L{0.2}|L{0.4}|L{0.4}|}
\hline
\textbf{Bảng} & \textbf{Chỉ mục} & \textbf{Mục đích} \\
\hline
\pkg{elections} & \pkg{(status, startDate, endDate)} & Lọc bầu cử theo trạng thái và thời gian \\
\hline
\pkg{election\_voters} & \pkg{unique(electionId, voterId)} & Tra cứu O(1) và loại bỏ phần tử trùng lặp \\
\hline
\pkg{votes} & \pkg{unique(electionId, voterId)} & Chống bỏ phiếu hai lần \\
 & \pkg{unique(electionId, blindedCommitment)} & Chống nộp commitment trùng \\
 & \pkg{(electionId, createdAt)} & Lấy theo thứ tự cho cây Merkle \\
\hline
\pkg{revealed\_votes} & \pkg{unique(electionId, revealKey)} & Chống phát lại reveal \\
 & \pkg{(electionId, candidateId)} & Tổng hợp kết quả theo ứng viên \\
\hline
\end{tabularx}
\end{table}

\textbf{Truy vấn song song:} Những tác vụ độc lập không phụ thuộc lẫn nhau được thực thi đồng thời bằng \pkg{Promise.all}. Ví dụ, khi hiển thị kết quả bầu cử, hệ thống đồng thời lấy dữ liệu từ MongoDB, Coordinator Service và Hyperledger Fabric thay vì chờ từng truy vấn hoàn thành tuần tự. Nhờ đó, thời gian phản hồi chỉ phụ thuộc vào truy vấn chậm nhất, giúp giảm đáng kể độ trễ khi tải dữ liệu.
